\documentclass[a4paper,12pt]{letter}
\usepackage[utf8]{inputenc}
\usepackage[english,slovene]{babel}
\usepackage[tikz]{bclogo}
\usepackage{qrcode,marvosym}
\usepackage{pgf,tikz,pgfplots}
\pgfplotsset{compat=1.14}
\usepackage{mathrsfs}
\usetikzlibrary{arrows}
\usepackage{graphicx}
\newcommand{\ds}{\displaystyle}
\usepackage{fancybox}
\usepackage{amsfonts}
\usepackage{amsmath}
\usepackage{wallpaper}
\usepackage{hyperref}
\usepackage[cp1250]{inputenc}
\usepackage{array}
\newcolumntype{M}[1]{>{\centering\arraybackslash}m{#1}}
\newcolumntype{N}{@{}m{0pt}@{}}
\usepackage{fancyhdr,enumerate}

\newcommand{\ora}{\overrightarrow}
\renewcommand{\headrulewidth}{3pt}
\renewcommand{\headrule}{\hbox to\headwidth{%
  \color{gray}\leaders\hrule height \headrulewidth\hfill}}
\renewcommand{\footrulewidth}{2pt}
\renewcommand{\footrule}{\hbox to\headwidth{%
  \color{brown}\leaders\hrule height \footrulewidth\hfill}}
\usepackage{gensymb}
\usepackage{amssymb}

\usepackage{fancyhdr}
\textwidth 20 true cm
\oddsidemargin -2 true cm
\evensidemargin -2 true cm
\topmargin -3.5 true cm
\textheight 27.5 true cm
\linespread{1.2}
\renewcommand{\headrulewidth}{0.3pt}
\renewcommand{\footrulewidth}{1.9pt}
\definecolor{qqwuqq}{rgb}{1,0,0}
\definecolor{qqqqff}{rgb}{0,0,1}
\definecolor{qqffqq}{rgb}{0,1,0}
\definecolor{ffqqqq}{rgb}{1,0,0}
\definecolor{zzuxux}{rgb}{0,0,0}
\usepackage{mdframed}
\usepackage{multicol}
\definecolor{bgblue}{RGB}{0,0,253}
\usepackage{xcolor}
\usepackage{eurosym}

\newcounter{tocke}
\setcounter{tocke}{0}
\usepackage{tikz-3dplot}

\mdfdefinestyle{theoremstyle}{%
linecolor=brown!40,linewidth=1pt,%
frametitlerule=true,%
frametitlebackgroundcolor=brown!50,
innertopmargin=\topskip,
}

\mdfdefinestyle{theoremstyle1}{%
linecolor=brown,middlelinewidth=1pt,%
frametitlerule=true,%
apptotikzsetting={\tikzset{mdfframetitlebackground/.append style={%
shade,left color=black!40, right color=gray!10},
mdfbackground/.append style={
top color=white!100!white,
bottom color=black!5
},
}},
frametitlerulecolor=gray,
frametitlerulewidth=2pt,
innertopmargin=\topskip,
innerbottommargin=\topskip,
}
\mdtheorem[style=theoremstyle1]{naloga}{Naloga}
\mdtheorem[style=theoremstyle]{resitev}{Analiza Naloge}
\mdtheorem[style=theoremstyle]{kriterij}{\v{S}tevilo dose\v{z}enih to\v{c}k na testu}

\usepackage[table]{xcolor}
\usepackage{titlesec}
\usepackage{venndiagram}
\usepackage[printwatermark]{xwatermark}
\usepackage{xcolor}
%\newwatermark[allpages,color=brown!2,angle=45,scale=5,
%xpos=-5,ypos=0]{matej.info}

\titleformat{\subsection}
  {\normalfont\Large\bfseries\color{brown}}
  {\thesubsection}{1em}{}

\titleformat{\subsubsection}
  {\normalfont\large\bfseries\color{brown}}
  {\thesubsubsection}{1em}{}

\begin{document}
\pagestyle{fancy}
\renewcommand\headrulewidth{0pt}
\lhead{}\chead{}\rhead{}
\rfoot{\vspace*{-2\baselineskip}}

\renewcommand{\vec}{\overrightarrow}
\everymath{\displaystyle}

\begin{minipage}{20cm}
\begin{bclogo}[couleur=brown!40,
  arrondi=0,ombre=true, couleurOmbre=gray!50,
  logo=\bcplume, noborder=true]
{\textrm{{\Large{Test 4.1} \Huge{}:
\large{Odvodi.} \hfill $PTI-2$}}}
\end{bclogo}
\end{minipage}

\begin{minipage}{20cm}
\begin{bclogo}[couleur=brown!40,ombre=true, couleurOmbre=gray!50,
  arrondi=0, logo=\bcplume, noborder=true]
% POPRAVEK: \setlength{\parskip}{} -> \setlength{\parskip}{0pt}
%\setlength{\parskip}{0pt}
{\textsc{{\large{Ime in priimek:}
\vspace{0.2cm} \rule{15cm}{0.2mm}\hfill }}}
\end{bclogo}
\end{minipage}


%==============================================
% NALOGE
%==============================================

\def\tocke2{4+3}
\begin{naloga}[\hfill \quad $\tocke2$
\qquad \rightsquigarrow \vert a.\qquad|b.\qquad|c.\qquad|]
\addtocounter{tocke}{\tocke2}
{Izračunaj odvod funkcije:}

a) $f(x)=5x^3-2x^2+7x$

b) $g(x)=\dfrac{5}{x^3}+1$
\end{naloga}
\vspace{7cm}

\def\tocke1{4+4}
\begin{naloga}[\hfill \quad $\tocke1$
\qquad \rightsquigarrow \vert a.\qquad|b.\qquad|c.\qquad|]
\addtocounter{tocke}{\tocke1}
{Izračunaj odvod v dani točki:}

a) $f(x)=(2x-1)(x-3)$, \quad $x_0=2$

b) $g(x)=\dfrac{3x+2}{x-2}$, \quad $x_0=4$
\end{naloga}
\newpage

\vspace{7cm}
\def\tocke3{6+3}
\begin{naloga}[\hfill \quad $\tocke3$
\qquad \rightsquigarrow \vert a.\qquad|b.\qquad|c.\qquad|]
\addtocounter{tocke}{\tocke3}
{Določi enačbo tangente na graf funkcije v dani točki:}

a) $f(x)=x^2$, \quad $x_0=2$. Nariši v koordinatni sistem.

b) $g(x)=\sqrt[3]{x^4}$, \quad $x_0=1$
\end{naloga}

\newpage
\def\tocke4{6}
\begin{naloga}[\hfill \quad $\tocke4$
\qquad \rightsquigarrow \vert a.\qquad|b.\qquad|c.\qquad|]
\addtocounter{tocke}{\tocke4}
{Izračunaj kot med krivuljama v presečišču:}
$f(x)=x^2-4$, \quad $g(x)=x^2-4x$
\end{naloga}

\newpage
\def\tocke5{6}
\begin{naloga}[\hfill \quad $\tocke5$
\qquad \rightsquigarrow \vert a.\qquad|b.\qquad|c.\qquad|]
\addtocounter{tocke}{\tocke5}
Pod kakšnim kotom seka graf funkcije $f(x)=\sqrt{x}-2$ os $x?$ Zapiši tudi presečišče.
\end{naloga}

\vfill
\mdtheorem[style=theoremstyle1]{kriterij}{Kriterij ocenjevanja}
\begin{kriterij*}[\hfill \v{s}tevilo vseh to\v{c}k na testu: $\arabic{tocke}$]
$$
\begin{array}{||c|c|c|c|c|c||>{}c|c|}
\hline\hline
\textup{ocena}& 1&2&3&4&5&\textrm{uspe\v{s}nost v } \%& \textrm{\bf OCENA}\\
\hline\hline
\% & [0,45) &[45,60)&[60,75) &[75,90) & [90,100]&
\mbox{\phantom{\huge{$\frac{5}{445}$15}}} & \\
\hline
\end{array}
\hspace*{0.5cm}\vspace*{-0.3cm}
\qrcode[hyperlink,height=0.8in]{http://www.matej.info}
$$
\vspace*{0.1cm}
\end{kriterij*}


%==============================================
% REŠITVE
%==============================================
\newpage
\begin{center}
{\Large\bfseries\color{brown} REŠITVE}
\end{center}

%----------------------------------------------
\begin{resitev}[Naloga 1 --- Odvod funkcije]

\textbf{a)} $f(x)=5x^3-2x^2+7x$

Uporabimo pravilo $(x^n)'=n\,x^{n-1}$:
$$
f'(x)=5\cdot3x^2 - 2\cdot2x + 7 = \boxed{15x^2-4x+7}
$$

\textbf{b)} $g(x)=\dfrac{5}{x^3}+1 = 5x^{-3}+1$
$$
g'(x)=5\cdot(-3)x^{-4} = \boxed{-\dfrac{15}{x^4}}
$$

\end{resitev}

\vspace{0.5cm}

%----------------------------------------------
\begin{resitev}[Naloga 2 --- Odvod v dani točki]

\textbf{a)} $f(x)=(2x-1)(x-3)$, \quad $x_0=2$

Najprej razpišemo produkt:
$$
f(x)=(2x-1)(x-3)=2x^2-6x-x+3=2x^2-7x+3
$$
$$
f'(x)=4x-7
$$
$$
f'(2)=4\cdot2-7=8-7=\boxed{1}
$$

\textbf{b)} $g(x)=\dfrac{3x+2}{x-2}$, \quad $x_0=4$

Kvocientno pravilo $\left(\dfrac{u}{v}\right)'=\dfrac{u'v-uv'}{v^2}$, kjer $u=3x+2$, $v=x-2$:
$$
g'(x)=\frac{3(x-2)-(3x+2)\cdot1}{(x-2)^2}=\frac{3x-6-3x-2}{(x-2)^2}=\frac{-8}{(x-2)^2}
$$
$$
g'(4)=\frac{-8}{(4-2)^2}=\frac{-8}{4}=\boxed{-2}
$$

\end{resitev}

\newpage

%----------------------------------------------
\begin{resitev}[Naloga 3 --- Enačba tangente]

Enačba tangente v točki $x_0$:\quad $t:\;y=f'(x_0)(x-x_0)+f(x_0)$

\textbf{a)} $f(x)=x^2$, \quad $x_0=2$
$$
f(2)=4,\quad f'(x)=2x,\quad f'(2)=4
$$
$$
t:\;y=4(x-2)+4=\boxed{4x-4}
$$

\begin{center}
\begin{tikzpicture}[scale=0.7]
  \begin{axis}[
    axis lines=center,
    xlabel={$x$}, ylabel={$y$},
    xmin=-1, xmax=4,
    ymin=-2, ymax=12,
    xtick={-1,0,...,4},
    ytick={0,2,...,12},
    grid=both,
    grid style={line width=0.2pt, draw=gray!30},
    tick label style={font=\small},
    width=9cm, height=8cm,
  ]
    \addplot[domain=-0.5:3.5, samples=100, thick, color=brown, line width=1.5pt]
      {x^2} node[above left, pos=0.85] {$f(x)=x^2$};
    \addplot[domain=0.5:3.5, samples=50, thick, color=blue, line width=1.5pt]
      {4*x-4} node[right, pos=0.9] {$t:y=4x-4$};
    \addplot[only marks, mark=*, mark size=2.5pt, color=red]
      coordinates {(2,4)};
    \node[above left, color=red] at (axis cs:2,4) {$(2,4)$};
  \end{axis}
\end{tikzpicture}
\end{center}

\textbf{b)} $g(x)=\sqrt[3]{x^4}=x^{4/3}$, \quad $x_0=1$
$$
g(1)=1,\quad g'(x)=\tfrac{4}{3}x^{1/3},\quad g'(1)=\tfrac{4}{3}
$$
$$
t:\;y=\tfrac{4}{3}(x-1)+1=\boxed{\tfrac{4}{3}x-\tfrac{1}{3}}
$$

\end{resitev}

\newpage

%----------------------------------------------
\begin{resitev}[Naloga 4 --- Kot med krivuljama]

$f(x)=x^2-4$, \quad $g(x)=x^2-4x$

\textbf{Presečišča:} $f(x)=g(x)$
$$
x^2-4=x^2-4x \;\Rightarrow\; 4x=4 \;\Rightarrow\; x=1
$$
Presečišče: $(1,\,-3)$

\textbf{Smerni koeficienti v presečišču:}
$$
f'(x)=2x \;\Rightarrow\; k_1=f'(1)=2
$$
$$
g'(x)=2x-4 \;\Rightarrow\; k_2=g'(1)=-2
$$

\textbf{Kot $\varphi$ med tangentama:}
$$
\tan\varphi=\left|\frac{k_1-k_2}{1+k_1 k_2}\right|
=\left|\frac{2-(-2)}{1+2\cdot(-2)}\right|
=\left|\frac{4}{-3}\right|=\frac{4}{3}
$$
$$
\varphi=\arctan\!\left(\tfrac{4}{3}\right)\approx\boxed{53{,}13^\circ}
$$

\end{resitev}

\vspace{0.5cm}

%----------------------------------------------
\begin{resitev}[Naloga 5 --- Kot s koordinatno osjo $x$ in presečišče]

$f(x)=\sqrt{x}-2$

\textbf{Presečišče z osjo $x$:} $f(x)=0$
$$
\sqrt{x}-2=0 \;\Rightarrow\; \sqrt{x}=2 \;\Rightarrow\; x=4
$$
$$
\boxed{\text{Presečišče: }(4,\,0)}
$$

\textbf{Smerni koeficient tangente v presečišču:}
$$
f'(x)=\frac{1}{2\sqrt{x}} \;\Rightarrow\; k=f'(4)=\frac{1}{2\cdot2}=\frac{1}{4}
$$

\textbf{Kot s pozitivno smerjo osi $x$:}
$$
\varphi=\arctan(k)=\arctan\!\left(\tfrac{1}{4}\right)\approx\boxed{14{,}04^\circ}
$$

\end{resitev}

\end{document}
